Ruang lingkup penelitian ini ditetapkan untuk menjaga fokus kajian serta memastikan proses kerja sama berjalan secara terarah dan realistis.
Bahasa isyarat yang digunakan dalam penelitian ini adalah Bahasa Isyarat Indonesia (BISINDO) dengan cakupan wilayah penggunaan terbatas pada komunitas di Jawa Barat.
Pembatasan wilayah dilakukan untuk menjaga konsistensi variasi isyarat serta memudahkan proses validasi linguistik bersama mitra.

Kelas isyarat yang digunakan dalam pembentukan dataset dibatasi pada tiga belas kata dasar BISINDO sebagaimana ditunjukkan pada Tabel~\ref{tab:kelas-isyarat-bisindo}.

\begin{table}[h]
\centering
\caption{Daftar kelas isyarat BISINDO yang digunakan}
\label{tab:kelas-isyarat-bisindo}
\begin{tabular}{ll}
\hline
No & Kelas isyarat \\
\hline
1 & \textit{teman} \\
2 & \textit{sahabat} \\
3 & \textit{maaf} \\
4 & \textit{izin} \\
5 & \textit{terima kasih} \\
6 & \textit{belum} \\
7 & \textit{tidak punya} \\
8 & \textit{hobi} \\
9 & \textit{hati-hati} \\
10 & \textit{ulang} \\
11 & \textit{halo} \\
12 & \textit{permisi} \\
13 & \textit{tolong} \\
\hline
\end{tabular}
\end{table}

Pemilihan kata dasar bertujuan untuk merepresentasikan variasi gerakan tangan, orientasi tubuh, serta ekspresi non-manual tanpa memperluas kompleksitas data secara berlebihan.

Penelitian ini menggunakan pendekatan spasial berbasis \textit{Residual Network} dua dimensi (ResNet-2D) sebagai model utama.
Data masukan model tidak berupa citra RGB mentah, melainkan representasi koordinat \textit{landmark} tubuh hasil ekstraksi MediaPipe yang terdiri atas nilai posisi dan perubahan posisi titik, yaitu $(x, y, \Delta x, \Delta y)$.
Pendekatan ini dipilih untuk meminimalkan ketergantungan pada informasi visual personal serta meningkatkan efisiensi pemodelan.

Penggunaan model \textit{Convolutional Neural Network} (CNN) standar diterapkan sebagai pembanding untuk memberikan konteks performa terhadap model ResNet-2D.
Data visual berupa citra atau video hanya digunakan pada tahap ekstraksi \textit{landmark} dan dokumentasi proses, serta tidak menjadi masukan langsung model.
Identitas visual partisipan dijaga melalui penerapan penyamaran area mata atau elemen identitas lain sesuai dengan prinsip perlindungan privasi.